\RequirePackage{plautopatch}
\RequirePackage[l2tabu, orthodox]{nag}

\documentclass[dvipdfmx]{jlreq}
\usepackage{graphicx}
\usepackage{pxrubrica}
\usepackage{enumitem}
\usepackage{url}
%\usepackage{wrapfig}
%\usepackage{floatflt}
\usepackage[dvipdfmx]{hyperref}
\usepackage{pxjahyper}
%\usepackage{alphalph}
%\renewcommand{\theenumi}{\AlphAlph{\value{enumi}}}
\title{\Huge 教授応募書類}

\author{\Large 鵜飼孝典}
\date{2026年5月27日}
\begin{document}

\large 国立大学法人北陸先端科学技術大学院大学御中

\vspace{50mm}
\begin{center}
{\Huge 教授応募書類 \par}%修士論文と記載される部分
\vspace{50mm}

{\Large 2026年5月27日 \par}% 氏名
\vspace{40mm}
{\Large 富士通株式会社 (シニアリサーチャー) \par }% 氏名
{\huge 鵜飼 孝典 \par}% 氏名
\end{center}
%\maketitle

\newpage
\begin{center}
\LARGE
履歴書
\end{center}
\includegraphics[keepaspectratio,width=35mm]{DSC_0326.png}

%\begin{floatingfigure}[r]{30mm}
% \includegraphics[keepaspectratio,width=30mm]
%  {DSC_0326.png}
%\end{floatingfigure}

\begin{description}[style=nextline]
\item [〇氏名（フリガナ）・性別] 鵜飼孝典（うかい たかのり）・男
\item [〇国籍] 日本
\item [〇生年月日(年齢)] 1966年7月8日(59歳)
\item [〇現住所] 〒182-0021\\
                 東京都調布市調布ヶ丘3-22-5
\item [〇連絡先] 自宅住所\\
                 Tel (042)485-3396\\
                 E-mail \href{mailto:takanori.ugai@gmail.com}{takanori.ugai@gmail.com}
\item [〇勤務先] 富士通株式会社 (シニアリサーチャー)\\
                 Tel (044)754-2652\\
                 E-mail \href{mailto:takanori.ugai@gmail.com}{takanori.ugai@gmail.com}
\item [〇学歴]
1985.3 　埼玉県立浦和高等学校 卒業\\
1986.4 　東京工業大学5類 入学\\
1990.3 　東京工業大学工学部計算工学科 卒業\\
1990.4 　東京工業大学電気電子工学専攻修士課程 入学\\
1992.3 　同上 修了\\
2009.4　 東京工業大学計算工学専攻博士課程 入学\\
2012.3　 同上 修了, 博士(工学)\\
　　「(論文題目) ステークホルダ分析とゴールグラフの品質を用いた要求獲得支援の研究」
\item [〇職歴]
1992.4～2021.3 　株式会社富士通研究所（2021.4 会社吸収に伴い移籍）\\
1996.7～1998.7 　英国 APM Limited 駐在\\
2005.7～2007.7 　米国 PARC社との協業プロジェクトに参加\\
2021.4～現在 　　富士通株式会社 シニアリサーチャー\\
2021.7～2026.3 　国立研究開発法人産業技術総合研究所
\item [〇専門分野] \underline{知識推論},知識獲得,知識共有,要求工学
\item [〇学会及び社会における活動等]
1992.4～  情報処理学会会員\\
2002.2～  ACM会員\\
人工知能学会会員
\end{description}

\newpage
\setcounter{section}{1}
\section{教育活動}

\begin{enumerate}
\item 担当講義・学内(社内)で実施した教育活動等\\
(科目等ごとに目標・内容と取り組んできた工夫等について記載)
\begin{enumerate}
\item 2006年7月 早稲田大学人間理工学部 非常勤講師\\
内容と工夫: 企業内における知識共有に関する講義。ソフトウェア開発企業における知識共有を自分の経験をもとに紹介した。ソフトウェア開発プロジェクト内で共有される知識、プロジェクト間で共有される知識、ソフトウェア開発に関するもの以外の共有される知識に分類して、それぞれを具体的に紹介し、社内で使われるツールを画像と共に紹介した。また組織のマネージャ、プロジェクトマネージャ、エンジニアのそれぞれの立場でツールを使う様子を画像を用いて紹介した。これにより知識共有が必要な理由に加え、立場により必要な知識が違うこと、同じ情報でも立場の違いによって知識としての見方が異なることを学生が納得できるようにした。
\item 2012年から2014年 筑波大学システム情報工学研究科 非常勤講師\\
担当科目「コンピュータサイエンス特別講義I」\\
最新のインターネット上のサービスに関する講義。SNSのサービスが日本でも始まったころに本講義を実施することになったため、一部の先進的なユーザの利用から、企業や政府機関による情報発信に続き、BLOGなどによる個人からの情報発信、SNSを用いた個人間コミュニケーションツールとしてのインターネットの変遷をそれぞれの盛衰を当時の画面例を示しながら説明した。同時にインターネットの技術が企業内のイントラネットでも用いられた流れを合わせて説明した。当時地域SNSなど大小さまざまな目的の異なるコミュニティを支援するサービスが立ち上がり始めたころであったため、この講義では、地域と大学を活性化することを目的としたSNSサービスのデザインを小グループで行うことを演習として実施した。この演習により様々な立場の人が利用するサービスについて、それぞれの立場からのメリット、セキュリティ、ユーザビリティなどについて、座学では納得しにくい課題について学生がそれぞれ気づいて検討してもらうことができた。
\item 2015年FUJITSUファミリ会 LS研究会 テクニカルアドバイザー\\
テーマ：オープンデータの活用可能性に関する研究　(Leading-edge Systems研究 最優秀賞受賞)\\
14の企業から一人づつ参加して半年間研究し、最後に派遣元の企業幹部が集まる場にて成果を発表する研究会。自分はこの研究会にテクニカルアドバイザーとして参加。参加メンバーは、全員オープンデータに興味はあるものの、どういうものかわからないもの、実際に使ってみたことがあるもの、インフラ企業で自社が持つデータをオープンデータとして公開することを検討しているものなど様々な立場から集まっていた。そこで、グループをオープンデータを使う立場と自社のデータをオープンデータとして公開して使ってもらう立場の２つに分けた。前半では、それぞれの立場で既存のデータ、ツール、技術の調査を行った。後半では、企業活動においてオープンデータを利用するメリット、デメリット、自社のデータを公開するメリット、デメリットを議論し、最後に２つのグループが検討した結果を総合して、将来の理想的なオープンデータが流通する状態を描いた。
\item 2016年自治体職員向けオープンデータ講座：多摩市、行方市、八王子市\\
2016年官民データ活用推進基本法において、国及び地方公共団体はオープンデータに取り組むことが義務付けられた。自治体からの依頼で、どのようなデータをどのような形で公開すればよいかということについて民間の観点から講習会を実施した。オープンデータを使った既存の商業サービスを紹介しつつ、前年のLS研究会の成果を活用して、オープンデータを使うメリット、デメリットを説明し、自治体が持つデータを公開することの意義を納得できるようにした。さらに具体的な使い方、加工の仕方を示すことでどのような形式で公開するとよいかという点で経産省が提示しているガイドラインが求めていることも納得しやすいようにした。
\item 2020年から2021年 筑波大学システム情報工学研究科 非常勤講師\\
担当科目「情報リテラシー(演習)」\\
本科目は、全学の新入生向けの基礎的な科目であり、Internetに流れる情報に関する信頼性、情報を取り扱う時の著作権に対する対応の仕方、Wordを使ったレポートの作成方法、Powerpointを使ったプレゼンテーションの作成方法、発表の仕方などを学ぶものである。2020年からコロナ感染症のため、学生はオンデマンドでビデオコンテンツで学習し、質問をTeamsを使ったTV会議で受け付ける形式で行われた。対面で演習を行う場合には、講義時間内に質問できるがオンデマンドでは、質問しづらいようであったため、提出されたレポートに適宜フィードバックを行い、何回でも再提出することを認めることで補うことにした。それぞれの課題について、多い学生で10回程度やり取りを行い、全体としては対面よりもむしろコミュニケーションの量は多くなったと考えている。
\item 2022年〜現在 青山学院大学大学院 非常勤講師\\
担当科目「データサイエンス特論」
\item 2023年〜現在 東北大学大学院 非常勤講師\\
担当科目「課題解決型PBL演習」
\end{enumerate}
\end{enumerate}

\newpage
\section{研究業績}
\begin{enumerate}
\item 学術論文
\begin{enumerate}
\item 査読付き国際論文誌
\begin{enumerate}
\item RDF-star2Vec: RDF-star Graph Embeddings for Data Mining\\
Shusaku Egami, \underline{Takanori Ugai}, Masateru Oota, Kyoumoto Matsushita, Takahiro Kawamura, Kouji Kozaki, Ken Fukuda\\
IEEE Access, vol. 11, pp. 142030-142042, 2023
\item Synthesizing Event-centric Knowledge Graphs of Daily Activities using Virtual Space\\
Shusaku Egami, \underline{Takanori Ugai}, Mikiko Oono, Koji Kitamura, Ken Fukuda\\
IEEE Access, vol. 11, pp. 23857-23873, 2023
\item Converging Semantic Knowledge and Deep Learning for Medical Coding\\
Nuria Garcia-Santa, Beatriz Miguel,  \underline{Takanori Ugai}\\
International Journal of Privacy and Health Information Management, vol. 7 pp32-52, 2019
\item Motivation Problems for Knowledge Sharing: Why People would not Share Knowledge?
Aoyama, Kouji, \underline{Takanori Ugai}, Akihiko Obata, and Hiroaki Harada\\
The International Journal of KNOWLEDGE, CULTURE \& CHANGE MANAGEMENT, vol. 8 issue.11, pp1-8, 2009
\item A mathematical model of knowledge transfer and case studies\\
\underline{Takanori Ugai}, Kouji Aoyama, Jun Arima\\
The International Journal of KNOWLEDGE, CULTURE \& CHANGE MANAGEMENT, vol. 7 pp9-17, 2007
\end{enumerate}
\item 査読付き国内論文誌
\begin{enumerate}
\item Wikidataを対象としたPointer Networkに基づく日本語エンティティリンキングモデルの構築と評価\\
澤村勇輝, 森田武史, 江上周作, \underline{鵜飼孝典}, 福田賢一郎\\
人工知能学会論文誌, 第39巻 第6号, p. C-O42\_1, 2024
\item Constructing and Evaluating a Japanese Entity Linking Model for Wikidata based on Pointer Network\\
Yuki Sawamura, Takeshi Morita, Shusaku Egami, \underline{Takanori Ugai}, Ken Fukuda\\
Transactions of the Japanese Society for Artificial Intelligence, vol. 39, no. 6, p. C-O42\_1, 2024
\item 属性つきゴールグラフにおけるゴールの品質特性\\
\underline{鵜飼孝典}, 林晋平, 佐伯元司\\
情報処理学会論文誌, 第55巻 第2号, pp893－908, 2014
\item 要求獲得におけるステークホルダの偏りと不足を検出する可視化ツール\\
\underline{鵜飼孝典}, 林晋平, 佐伯元司\\
情報処理学会論文誌, 第53巻 第4号, pp1448－1460, 2012
\end{enumerate}
\item 査読付き国際会議
\begin{enumerate}[label=\arabic*,align=right, labelsep=3pt]
\item RAG-Enhanced Prompt Compression with Need-Oriented Knowledge for Dialog Based Embodied Navigation\\
Hiroaki Shimoma, Sudesna Chakraborty, Takeshi Morita, Aoi Ohta, Masaki Asada, Shusaku Egami, \underline{Takanori Ugai}, Masahiro Hamasaki\\
Proceedings of 18th International Conference on Agents and Artificial Intelligence (ICAART2026), 2026
\item Comparison of Metadata Representation Models for Knowledge Graph Embeddings\\
Shusaku Egami, Kyoumoto Matsushita, \underline{Takanori Ugai}, Ken Fukuda\\
Proceedings of Knowledge Graphs and Semantic Web Conference 2025 (KGSWC2025), 2025
\item Household Task Planning with Multi-Objects State and Relationship Using Large Language Models Based Preconditions Verification\\
Jin Aoyama, Sudesna Chakraborty, Takeshi Morita, Shusaku Egami, \underline{Takanori Ugai}, Ken Fukuda\\
Proceedings of the 17th International Conference on Agents and Artificial Intelligence (ICAART2025), pp. 472-483, 2025
\item Reasoning and Justification System for Domestic Hazardous Behaviors Based on Knowledge Graph of Daily Activities and Retrieval-Augmented Generation\\
Fumikatsu Anaguchi, Sudesna Chakraborty, Takeshi Morita, Shusaku Egami, \underline{Takanori Ugai}, Ken Fukuda\\
Proceedings of the Twelfth International Symposium on Computing and Networking (CANDAR2024), 2024
\item Entity Linking for Wikidata using Large Language Models and Wikipedia Links\\
Fumiya Mitsuji, Sudesna Chakraborty, Takeshi Morita, Shusaku Egami, \underline{Takanori Ugai}, Ken Fukuda\\
Proceedings of the 9th International Workshop on GPU Computing and AI in conjunction with Twelfth International Symposium on Computing and Networking (CANDAR2024), 2024
\item Compressing Multi-Modal Temporal Knowledge Graphs of Videos\\
Shusaku Egami, \underline{Takanori Ugai}, Ken Fukuda\\
Proceedings of the ISWC 2024 Posters, Demos and Industry Tracks co-located with 23rd International Semantic Web Conference (ISWC2024), 2024
\item VHAKG: A Multi-modal Knowledge Graph Based on Synchronized Multi-view Videos of Daily Activities\\
Shusaku Egami, \underline{Takanori Ugai}, Swe Nwe, Nwe Htun, Ken Fukuda\\
Proceedings of the 33rd ACM International Conference on Information and Knowledge Management (CIKM2024), 2024
\item Multimodal Datasets and Benchmarks for Reasoning about Dynamic Spatio-Temporality in Everyday Environments\\
\underline{Takanori Ugai}, Kensho Hara, Shusaku Egami, Ken Fukuda\\
Embodied AI 2024 workshop in conjunction with The IEEE/CVF Conference on Computer Vision and Pattern Recognition 2024 (EAI-CVPR2024), 2024
\item Exploring Spatial Relation Awareness through Virtual Indoor\\
Swe Nwe Nwe Htun, Shusaku Egami, \underline{Takanori Ugai}, Yijun Duan, Ken Fukuda\\
Proceedings of the 26th International Conference on Human-Computer Interaction (HCII2024), pp. 34-51, 2024
\item Zero-Shot Query Experiments in Knowledge Graph Reasoning Challenge for Older Adults Safety\\
Ken Fukuda, \underline{Takanori Ugai}, Shusaku Egami, Kyoumoto Matsushita\\
2024 IEEE 18th International Conference on Semantic Computing (ICSC), pp. 301-305, 2024
\item Event Prediction in Event-Centric Knowledge Graphs Using BERT\\
\underline{Takanori Ugai}, Shusaku Egami, Ken Fukuda\\
2024 IEEE 18th International Conference on Semantic Computing (ICSC), pp. 306-310, 2024
\item Japanese Pointer Network based Entity Linker for Wikidata\\
Yuki Sawamura, Takeshi Morita, Shusaku Egami, \underline{Takanori Ugai}, Kenichiro Fukuda\\
Proceedings of The 12th International Joint Conference on Knowledge Graphs (IJCKG2023), 2023
\item Automatic Action Script Generation to Improve Execution Rate based on LLM in VirtualHome\\
Jin Aoyama, Takeshi Morita, \underline{Takanori Ugai}, Shusaku Egami, Ken Fukuda\\
Proceedings of The 12th International Joint Conference on Knowledge Graphs (IJCKG2023), 2023
\item Analysis of Annotation Quality of Human Activities Using Knowledge Graphs\\
Shusaku Egami, Mikiko Oono, Mai Otsuki, \underline{Takanori Ugai}, Ken Fukuda\\
Communications in Computer and Information Science, pp. 483-489, 2023
\item Datasets of Mystery Stories for Knowledge Graph Reasoning Challenge\\
Kozaki Kouji, Shusaku Egami, Kyoumoto Matsushita, \underline{Takanori Ugai}, Takahiro Kawamura, Ken Fukuda\\
Proceedings of the Workshop on Semantic Methods for Events and Stories (SEMMES) co-located with the 20th Extended Semantic Web Conference (ESWC2023), 2023
\item Criminal Investigation with Augmented Ontology and Link Prediction\\
\underline{Takanori Ugai}\\
2023 IEEE 17th International Conference on Semantic Computing (ICSC), pp. 288-289, 2023
\item Contextualized Scene Knowledge Graphs for XAI Benchmarking\\
Takahiro Kawamura, Shusaku Egami, Kyoumoto Matsushita, \underline{Takanori Ugai}, Ken Fukuda, Kouji Kozaki\\
Proceedings of the 11th International Joint Conference on Knowledge Graphs (IJCKG2022), pp. 64-72, 2022
\item The Magic of Semantic Enrichment and NLP for Medical Coding\\
Nuria García-Santa, Beatriz San Miguel, \underline{Takanori Ugai}\\
Extended Semantic Web Conference 2019 Satellite Events, pp58-63, 2019
\item Report on the First Knowledge Graph Reasoning Challenge 2018 - Toward the eXplainable AI System\\
Takahiro Kawamura, Shusaku Egami, Koutarou Tamura, Yasunori Hokazono, \underline{Takanori Ugai}, Yusuke Koyanagi, Fumihito Nishino, Seiji Okajima, Katsuhiko Murakami, Kunihiko Takamatsu, Aoi Sugiura, Shun Shiramatsu, Shawn Zhang, Kouji Kozaki\\
9th Joint International Conference, LNCS Vol. 12032 pp18-34, 2019
\item Conversion of Physical Quantity and its Application\\
\underline{Takanori Ugai}, Shohei Yamane\\
Proceedings of the International Semantic Web Conference 2017 Posters {\&} Demonstrations and Industry Tracks, \url{http://ceur-ws.org/Vol-1963/paper534.pdf} 2017
\item One concept expressed in five ways: How does THe Media format affect feedback?\\
Chie Sasaki, Keiko Ihara, \underline{Takanori Ugai}\\
Ethographyic Praxis in Industry Conference Proceedings 2013, p409 2013
\item Co-reniewal of a village with local residents: A case study of a fisherman's village\\
\underline{Takanori Ugai}\\
Ethographyic Praxis in Industry Conference Proceedings 2012, pp376-377 2012
\item What an Earthquake Changed\\
\underline{Takanori Ugai}\\
Ethographyic Praxis in Industry Conference Proceedings 2011, pp169-172 2011
\item Toward industrialization of ethnograpy\\
\underline{Takanori Ugai}, Kouji Aoyama, Akihiko Obata\\
Ethographyic Praxis in Industry Conference Proceedings 2010, pp26-35 2010
\item Virtual Ethnography System: analysis of word-of-Moouth on Blogsphere\\
Kouji Aoyama, Tetsuro Takahashi, \underline{Takanori Ugai}\\
Ethographyic Praxis in Industry Conference Proceedings 2010, ppp318-319 2010
\item Visualizing Stakeholder Concerns with Anchored Map\\
\underline{Takanori Ugai}, Kouji Aoyama\\
Human Interface and the Management of Information. Designing Information Environments, Symposium on Human Interface 2009 LNCS vol. 5617 pp181-189 2009
\item Organization Diagnosis Tools Based on Social Network Analysis\\
\underline{Takanori Ugai}, Kouji Aoyama\\
Human Interface and the Management of Information. Designing Information Environments pp181-189 2009
\item Domain Knowledge Wiki for Eliciting Requirements\\
\underline{Takanori Ugai}, Kouji Aoyama\\
2009 Second International Workshop on Managing Requirements Knowledge pp 4-6, 2009
\item Visualisation Tool for Cooperative Relations and Its Verification with a Case Study\\
\underline{Takanori Ugai}, Kouji Aoyama\\
12th International Conference on Knowledge-Base and Intelligent Information \& Engineering Systems LNCS vol.5178 pp.814-822 2008
\item Extraction of Viewpoints for Eliciting Customer's Requirements based on Analysis of Specification Change Records\\
Kouji Aoyama, \underline{Takanori Ugai}, Shigeru Yamada, Akihiko Obata\\
14th Asia-Pacific Software Engineering Conference pp33-40 2007
\item Design and Evaluation a Knowledge Management System by Using Mathematical Model of Knowledge Transfer\\
Kouji Aoyama, \underline{Takanori Ugai}, Jun Arima\\
Knowledge-Based Intelligent Information and Engineering Systems, 11th
International Conference, LNCS vol.4694 pp1253-1260 2007
\item An Interface of Multi-viewpoint Maps for Re-using Software Products\\
\underline{Takanori Ugai}, Terunobu Kume, Kazuo Misue\\
Proceedings of HCI International 2005, 2005-07
\item Maintenance Support of Corporate Directories with Social-filtering
Kazuo Misue, \underline{Takanori Ugai}\\
Proceedings of HCI International 2003, Human - Centred Computing, Cognitive, Social and Ergonomic Aspects/3/p.1061-1065, 2003-01
\item Link-based Acquisition of Web Metadata for Domain-specific Directories\\
Hiroshi Tsuda, \underline{Takanori Ugai}, Kazuo Misue\\
The 2000 Pacific Rim Knowledge Acquisition Workshop pp.317-324 2000
\item Security Policy Objects and RM-ODP\\
\underline{Takanori Ugai}\\
OMA/ODP Workshop 1997
\item Reflective Specification: Applying A Reflective Language to Formal Specification\\
Motoshi Saeki, Takeshi Hiroi,\underline{Takanori Ugai}\\
Proceedings of the 7th International Workshop on Software Specification and Design pp204-213 1993
\end{enumerate}
\item その他論文・研究発表（査読なし）
\begin{enumerate}[label=\arabic*,align=right, labelsep=3pt]
\item 欲求指向知識と RAG を用いた 対話ナビゲーションシステムのためのプロンプト圧縮手法\\
下間宇晟, チャクラボルティ シュデシナ, 森田武史, 太田葵, 浅田真生, 江上周作, \underline{鵜飼孝典}, 濱崎雅弘\\
2026年度人工知能学会全国大会（第40回）, 2026年6月
\item 大規模視覚言語モデルを用いた日常生活動画からの イベント中心知識グラフ生成\\
後藤颯志, チャクラボルティ シュデシナ, 森田武史, 太田葵, Imrattanatrai Wiradee, 浅田真生, 江上周作, \underline{鵜飼孝典}, 濱崎雅弘\\
言語処理学会第32回年次大会（NLP2026）, 2026年3月
\item コストと精度を考慮した質問と文書に基づく適応的 RAG 選択\\
穴口史将, チャクラボルティ シュデシナ, 森田武史, 太田葵, 浅田真生, 江上周作, \underline{鵜飼孝典}, 濱崎雅弘\\
言語処理学会第32回年次大会（NLP2026）, 2026年3月
\item 常識知識グラフを活用した動作ラベル間の等価関係の自動構築手法\\
三辻史哉, チャクラボルティ シュデシナ, 森田武史, 吉川友也, 山本泰智, 太田葵, 江上周作, \underline{鵜飼孝典}, 濱崎雅弘\\
言語処理学会第32回年次大会（NLP2026）, 2026年3月
\item 文章からの因果関係抽出用カテゴリ付きベンチマークデータセット\\
\underline{鵜飼孝典}\\
第67回セマンティックウェブとオントロジー研究会 2025(SWO-067), 2025年11月
\item 因果ナレッジグラフ生成に向けた AI マルチエージェントと大規模言語モデルを用いた因果関係抽出\\
\underline{鵜飼孝典}\\
第66回人工知能学会セマンティックウェブとオントロジー研究会 2025(SWO-066), 2025年9月
\item 大規模視覚言語モデルを用いた日常生活動画からのイベント中心知識グラフ生成の検討\\
後藤颯志, チャクラボルティ シュデシナ, 森田武史, 太田葵, Imrattanatrai Wiradee, 浅田真生, 江上周作, 濱崎雅弘, \underline{鵜飼孝典}\\
第66回人工知能学会セマンティックウェブとオントロジー研究会 2025(SWO-066), 2025年9月
\item 常識知識グラフを用いたメタ動作認識データセットの自動構築\\
三辻史哉, チャクラボルティ シュデシナ, 森田武史, 江上周作, 吉川友也, 山本泰智, \underline{鵜飼孝典}, 福田賢一郎\\
2025年度人工知能学会全国大会（第39回）, 2025年6月
\item 動画キャプション生成を用いた動作認識データセットの拡張\\
後藤颯志, チャクラボルティ シュデシナ, 森田武史, 江上周作, 吉川友也, 山本泰智, \underline{鵜飼孝典}, 福田賢一郎\\
2025年度人工知能学会全国大会（第39回）, 2025年6月
\item 文書のチャンクに基づくGraphRAGと知識グラフのコミュニティ分類に基づくGraphRAGの比較分析\\
穴口史将, チャクラボルティ シュデシナ, 森田武史, 江上周作, \underline{鵜飼孝典}, 福田賢一郎\\
2025年度人工知能学会全国大会（第39回）, 2025年6月
\item Embedding Method for Knowledge Graph with Densely Defined Ontology\\
\underline{Takanori Ugai}\\
arXiv, 2025年4月
\item 動画キャプション生成とMetaVDを用いた動作認識データセットの拡張\\
後藤颯志, 森田武史, チャクラボルティ シュデシナ, 吉川友也, 山本泰智, 江上周作, \underline{鵜飼孝典}, 福田賢一郎\\
第65回人工知能学会セマンティックウェブとオントロジー研究会, 2025年3月
\item 大規模言語モデルを用いたオブジェクトの状態と関係に基づく家庭内タスクプランニング\\
青山仁, チャクラボルティ シュデシナ, 森田武史, \underline{鵜飼孝典}, 江上周作, 福田賢一郎\\
第63回人工知能学会セマンティックウェブとオントロジー研究会, 2024年9月
\item 3D Spatial Grounding in Virtual Environments: A VirtualHome2KG Study\\
Swe Nwe Nwe Htun, Shusaku Egami, \underline{Takanori Ugai}, Yijun Duan, Ken Fukuda\\
第27回画像の認識・理解シンポジウム（MIRU2024）, 2024年6月
\item LLMの常識知識を活用した日常生活データセット自動構築手法の提案\\
青山仁, 森田武史, \underline{鵜飼孝典}, 江上周作, 福田賢一郎\\
2024年度人工知能学会全国大会論文集（第38回）, 2024年6月
\item 第2回国際ナレッジグラフ推論チャレンジ：日常生活に関するマルチモーダルデータからの行動の予測に対するLLMへの適用\\
\underline{鵜飼孝典}, 江上周作, 川村隆浩, 古崎晃司, 森田武史, 松下京群, 小川智広, 吉岡寛悟, 平野司, 尾﨑健吾, 福田賢一郎\\
2024年度人工知能学会全国大会論文集（第38回）, 2024年5月
\item DBpediaオントロジーとGPTに基づくWikipediaの赤リンクを用いたDBpediaの拡張\\
森俊人, 森田武史, \underline{鵜飼孝典}, 江上周作, 福田賢一郎\\
セマンティックウェブとオントロジー研究会 2023(SWO-062), 2024年3月
\item マルチモーダル大規模言語モデルと画像キャプションに基づく描画内容に即した併置型駄洒落の認識\\
浅野歴, 森田武史, \underline{鵜飼孝典}, 江上周作, 福田賢一郎\\
セマンティックウェブとオントロジー研究会 2023(SWO-062), 2024年3月
\item ナレッジグラフ推論チャレンジ2023開催報告-生成AI時代のナレッジグラフ構築技術\\
古崎晃司, 江上周作, 松下京群, \underline{鵜飼孝典}, 川村隆浩, 福田賢一郎\\
セマンティックウェブとオントロジー研究会 2023(SWO-062), 2024年3月
\item LLMを利用したイベント中心ナレッジグラフにおけるイベント予測\\
\underline{鵜飼孝典}\\
第61回セマンティックウェブとオントロジー研究会 2023(SWO-061), 2023年11月
\item LLMを活用した抽象的なタスク記述からのVirtualHomeのためのアクションスクリプト自動生成\\
青山仁, 森田武史, \underline{鵜飼孝典}, 江上周作, 福田賢一郎\\
第61回セマンティックウェブとオントロジー研究会 2023(SWO-061), 2023年11月
\item 述語の意味によるクラスタリングを用いたシーングラフ生成\\
太田雅輝, \underline{鵜飼孝典}, 江上周作, 清雄一, 田原康之, 大須賀昭彦, 福田賢一郎\\
第60回人工知能学会セマンティックウェブとオントロジー研究会 SIG-SWO-060-02, 2023年8月
\item ナレッジグラフ推論チャレンジ【実社会版】開催報告\\
\underline{鵜飼孝典}, 江上周作, 川村隆浩, 松下京群, 古崎晃司, 福田賢一郎\\
第60回人工知能学会セマンティックウェブとオントロジー研究会 SIG-SWO-060-03, 2023年8月
\item 大規模知識グラフを対象とした英語エンティティリンキングモデルの日本語対応における課題の分析\\
澤村勇輝, 谷津元樹, 森田武史, 江上周作, \underline{鵜飼孝典}, 福田賢一郎\\
2023年度人工知能学会全国大会論文集, 2023年6月
\item 第1回国際版ナレッジグラフ推論チャレンジ2023開催報告\\
古崎晃司, 江上周作, 松下京群, \underline{鵜飼孝典}, 川村隆浩, 福田賢一郎\\
2023年度人工知能学会全国大会論文集, 2023年6月
\item 家庭内の日常生活動画とイベント中心知識グラフの同時生成 ― ナレッジグラフ推論チャレンジ【実社会版】の開催に向けて ―\\
江上周作, \underline{鵜飼孝典}, Swe Nwe Nwe Htun, 太田雅輝, 大野美喜子, 北村光司, 松下京群, 古崎晃司, 川村隆浩, 福田賢一郎\\
2023年度人工知能学会全国大会論文集, 2023年6月
\item 画像キャプション生成及び物体認識を用いた駄洒落文ランキング手法の評価\\
浅野歴, 谷津元樹, 森田武史, 江上周作, \underline{鵜飼孝典}, 福田賢一郎\\
言語処理学会第29回年次大会発表論文集, 2023年3月
\item シーングラフ生成におけるロングテール問題解決に向けたデータサンプリング手法の検討\\
太田雅輝, \underline{鵜飼孝典}, 江上周作, 清雄一, 田原康之, 大須賀昭彦, 福田賢一郎\\
第59回セマンティックウェブとオントロジー研究会 2023(SWO-059), 2023年3月
\item Wikipediaの赤リンクを用いたDBpediaの拡張の検討\\
森俊人, 谷津元樹, 森田武史, 江上周作, \underline{鵜飼孝典}, 福田賢一郎\\
第59回セマンティックウェブとオントロジー研究会 2023(SWO-059), 2023年3月
\item イベント中心ナレッジグラフにおけるWalkモデルと代数モデルの埋め込みの比較\\
\underline{鵜飼孝典}\\
第59回セマンティックウェブとオントロジー研究会 2023(SWO-059), 2023年3月
\item Wikidataを対象とした英語エンティティリンキングモデルの日本語対応における課題の分析\\
澤村勇輝, 谷津元樹, 森田武史, 江上周作, \underline{鵜飼孝典}, 福田賢一郎\\
第59回セマンティックウェブとオントロジー研究会 2023(SWO-059), 2023年3月
\item シーングラフ生成の精度向上に向けた最適なデータセット生成の調査\\
太田雅輝, 江上周作, \underline{鵜飼孝典}, 福田賢一郎\\
第58回セマンティックウェブとオントロジー研究会 2022(SWO-058), 2022年11月
\item イベント中心ナレッジグラフにおける埋め込みを用いたイベントのタイプ予測\\
\underline{鵜飼孝典}\\
第58回セマンティックウェブとオントロジー研究会 2022(SWO-058), 2022年11月
\item 動画とナレッジグラフを併用した日常生活の表現を支援しリスクを可視化するツール\\
\underline{鵜飼孝典}, 江上周作, 窪田文也, 大野美喜子, 福田賢一郎\\
第58回セマンティックウェブとオントロジー研究会 2022(SWO-058), 2022年11月
\item Daily Activity data generation in Cyberspace for Semantic AI technology and HRI simulation\\
Ken Fukuda, \underline{Takanori Ugai}, Mikiko Oono, Koji Kitamura, Fumiya Kubota, Takeshi Morita, Swe Nwe Nwe Htun, Shusaku Egami\\
第40回日本ロボット学会学術講演会, 2022年9月
\item イベント中心ナレッジグラフにおけるリンク予測の予測モデルによる違い\\
\underline{鵜飼孝典}, 江上周作, 福田賢一郎\\
第57回人工知能学会セマンティックウェブとオントロジー研究会 2022(SWO-057), 2022年8月
\item イベント中心ナレッジグラフ埋め込みにおけるメタデータ表現モデルの分析\\
江上周作, \underline{鵜飼孝典}, 太田雅輝, 川村隆浩, 松下京群, 古崎晃司\\
第57回人工知能学会セマンティックウェブとオントロジー研究会 2022(SWO-057), 2022年6月
\item コンペティションによる協創：安心安全を守るAIの開発に向けて\\
\underline{鵜飼孝典}, 江上周作, 大野美喜子, 福田賢一郎, 川村隆浩, 古崎晃司, 松下京群\\
第191回ヒューマンインタフェース学会研究会「社会のデザイン・市民のデザイン(SIG-UXSD-15）」, 2022年6月
\item イベント中心知識グラフによる人間生活を含む環境のサイバー空間への転写にむけて\\
福田賢一郎, 江上周作, \underline{鵜飼孝典}, 森田武史, 大野美喜子, 北村光司, Qiu Yue, 原健翔, 古崎晃司, 川村隆浩\\
2022年度人工知能学会全国大会論文集, 2022年6月
\item 説明生成のための知識グラフ構築ガイドラインの考察\\
古崎晃司, 江上周作, 松下京群, \underline{鵜飼孝典}, 川村隆浩\\
人工知能学会全国大会論文集 第36回 (2022), 2022年6月
\item 高齢者の家庭内事故予防に役立つAIシステムの開発 —産業版ナレッジグラフ推論チャレンジに向けて—\\
\underline{鵜飼孝典}, 江上周作, 大野美喜子, 窪田文也, 福田賢一郎, 川村隆浩, 古崎晃司, 松下京群\\
人工知能学会第二種研究会資料 2022(SWO-056), 2022年3月
\item 家庭内の事故予防に向けた合成ナレッジグラフの構築と推論\\
江上周作, \underline{鵜飼孝典}, 窪田文也, 大野美喜子, 北村光司, 福田賢一郎\\
人工知能学会第二種研究会資料 2022(SWO-056), 2022年3月
\end{enumerate}
\end{enumerate}
\end{enumerate}

\item 主要論文・著書の概要
\begin{itemize}
\item 1-(a)-i: Converging Semantic Knowledge and Deep Learning for Medical Coding\\
カルテの記述に病名を付けることは保険請求、公的統計などの目的で使われる重要なタスクである。
本論文では、病名の語彙集（オントロジー）を医学論文を用いて拡張し、
既存のテキスト距離、テキスト埋め込みなどを組み合わせた手法により
このタスクの精度を上げた。提案手法では、学習データを必要とせず、新しいコーディングシステムや
国や精度によって異なるコーディングシステムに対しても有効である。また、病名をオントロジーと
医学論文を組み合わせて拡張しているため、カルテ中の短い記述に対しても適切な病名を
つけることができた。

\item 1-(a)-iii: A mathematical model of knowledge transfer and case studies\\
企業内の知識共有、知識移転を数学的に表現するモデルを提案した論文。
このモデルは, 知識の流れを提供側と受取側に分けて,
それぞれの利益, コスト, 提供側と受取側の間に存在する障壁を
考えることで, プロジェクト内の協力関係を知識移転の
観点から分析することができる.
利益からコストを差し引いて，なおかつ障壁よりも多くの利益が得られ
る場合に知識移転が進むと考える.
このモデルは, 既存のナレッジマネジメント施策の成功/失敗の
原因分析, 新しい意識共有システム設計のための
フィールド調査と設計したシステムの評価などに活用されている.

\item 1-(b)-i : 属性つきゴールグラフにおけるゴールの品質特性\\
ソフトウェア開発の要求分析のフェーズで用いるゴール指向要求分析法で得られるゴールグラフの
各ゴールに対する品質特性
を定義し, ゴールグラフの品質を下げる要因になるゴールを同定できるように
する技術を提案した論文。
それぞれの品質特性に対して, 各ゴールが満たすべき必要条件を
属性付きゴールグラフの属性を用いて品質特性述語として定義した.
例えば、ゴールに対するステークホルダの理解, 解釈が一致していない場合、このゴール
は, 非あいまい性を満たしていないと定義した.
定義したゴールに対する品質特性と品質特性述語が品質的に問題のあるゴールの発見と修正を効率的に
行うために役立つ.
\item 1-(b)-ii: 要求獲得におけるステークホルダの偏りと不足を検出する可視化ツール\\
高品質な要求獲得を行うためには, 各要件に対して, 適切なステークホルダによって検討されることが重要となる.
ステークホルダと要件の関係を可視化した提案ツールでは,
枝で結ばれたノード間には, 距離が離れれば離れる程, また枝に付与された値が
大きければ大きい程大きな引力が働くので,ステークホルダが, その要件に強い関心があるとき, お互いにより
近く配置される. このように表現することによって, 強い関心を持つステークホルダが
いない関心事の近くにはステークホルダが配置されず,
また誰も関心を持っていない場合にはどのステークホルダもその要件に
つながらない. これにより,プロジェクトマネージャや
分析者が関心事に対するステークホルダの偏りや不足を知ることができる.

\item 1-(c)-8: Multimodal Datasets and Benchmarks for Reasoning about Dynamic Spatio-Temporality in Everyday Environments\\
家庭環境における動的かつ時空間的な出来事に対する推論を行うEmbodied AI（身体性人工知能）の発展を目的とし、3Dシミュレータを用いてアノテーション済みの日常生活動画データを生成する手法とデータセットを提案した。具体的には、動画から人間の日常生活や周囲の環境に対するAIの理解度を測定するためのマルチモーダルQAデータセット「MMQADL (Multimodal Question Answering Dataset of Daily Life)」を構築した。このデータセットには、場所、動作、対象物、時間などに関する記述的および定量的な質問が含まれており、Video-LLaVaやGemini 1.5 Proなどの画像・映像生成AIモデルを用いて評価実験を行い、ロボット等の自律システムにおける日常生活理解のベンチマークを示した。

\item 1-(c)-24: Toward industrialization of ethnograpy\\
文化人類学や社会学で行われていた行動観察調査には、現場の内情や考え方など、現場のリアルな様子を理解できるという
メリットがある。そのため、企業や顧客の課題解決、あるいは課題発見のために用いられることが多くなった。
しかしながら, 人員、コスト、品質において見積もりが難しいことがビジネスとして行う上での課題の一つになっている。
本論文では、ソフトウェア開発プロジェクトのプロセスを改善するために定義された能力成熟度モデル統合(CMMI)を
行動観察調査に当てはめて、定義しなおした。成果物は、プロジェクトや構成員の品質だけではなく、
観察対象になる現場の状況にも依存してしまうため、このモデルでは、プロジェクトの管理、プロセスを
評価する。このモデルにより顧客にたいして行動観察調査を用いたプロジェクトの見積もりについて
よりはっきりした根拠を示すことができるようになった。
\end{itemize}
\item 特許
\begin{enumerate}[label=\arabic*,align=parleft]
\item 出願番号JP2021/009344  機械学習プログラム、予測プログラム、装置、及び方法 \underline{鵜飼孝典} 2021
\item 出願番号JP2020/041077  機械学習プログラム、推定プログラム、装置、及び方法 \underline{鵜飼孝典} 2020
\item 出願番号JP2020/016804  オントロジー生成プログラム、オントロジー生成装置およびオントロジー生成方法 \underline{鵜飼孝典} 2020
\item 出願番号JP2019/046458  学習方法、学習装置および学習プログラム 岡嶋成司,\underline{鵜飼孝典} 2019
\item 公開番号WO2021/064879 学習方法、学習装置、学習プログラム、予測方法、予測装置および予測プログラム \underline{鵜飼孝典},岡嶋成司 2019
\item 登録番号6798307 変換プログラム、変換方法、及び変換装置 \underline{鵜飼孝典} 2016
\item 登録番号10353667 情報変換方法、情報処理装置、及び情報変換プログラム \underline{鵜飼孝典},山根昇平 2015
\item 登録番号6657764 情報変換方法、情報処理装置、及び情報変換プログラム \underline{鵜飼孝典},山根昇平 2014
\item 登録番号10360243 情報提示プログラム、情報提示方法及び情報提示装置 \underline{鵜飼孝典} 2013
\item 登録番号6557959 情報提示プログラム、情報提示方法及び情報提示装置 \underline{鵜飼孝典} 2012
\item 公開番号3016006 情報提示プログラム、情報提示方法及び情報提示装置 \underline{鵜飼孝典} 2011
\item 登録番号5578091 人間関係マップ作成支援装置，人間関係マップ作成支援方法および人間関係マップ作成支援プログラム \underline{鵜飼孝典} 2010
\item 登録番号5423293 グラフ表示プログラムおよびグラフ表示方法 \underline{鵜飼孝典} 2010
\item 登録番号5240024 組織管理支援装置 \underline{青山浩二},小林紀之,鵜飼孝典,原田裕明 2010
\item 登録番号5293415 評価装置、評価方法および評価プログラム \underline{鵜飼孝典} 2009
\item 登録番号5343642 活性度測定プログラム \underline{鵜飼孝典},青山浩二,小林紀之 2009
\item 公開番号2009-0006381 情報検索装置、情報検索方法および情報検索プログラム 青山浩二,小林　紀之,\underline{鵜飼孝典},有馬淳 2009
\item 登録番号4978336 情報検索装置、情報検索方法および情報検索プログラム 青山浩二,小林紀之,\underline{鵜飼孝典},有馬淳 2009
\item 公開番号2008-0028362 振り返りデータ処理方法、振り返りデータ評価方法及び装置 \underline{鵜飼孝典},青山浩二,有馬淳,小林紀之 2008
\item 登録番号5130732 振り返りデータ処理方法、振り返りデータ評価方法及び装置 \underline{鵜飼孝典},青山浩二,有馬淳,小林紀之 2008
\item 公開番号2008-0120623 業務フロー管理プログラム、業務フロー管理装置、および業務フロー管理方法 原田裕明,有馬淳,小幡明彦,\underline{鵜飼孝典} 2008
\item 登録番号4984846 業務フロー管理プログラム、業務フロー管理装置、および業務フロー管理方法 原田裕明,有馬淳,小幡明彦,\underline{鵜飼孝典} 2008
\item 公開番号2008-0077421 施策管理プログラム、施策管理装置、施策管理方法 小林紀之,有馬淳,\underline{鵜飼孝典},青山浩二 2008
\item 登録番号4957145 施策管理プログラム、施策管理装置、施策管理方法 小林紀之,有馬淳,\underline{鵜飼孝典},青山浩二 2008
\item 公開番号2008-0091450 改善活動支援プログラム、方法及び装置 有馬淳,\underline{鵜飼孝典},青山浩二,小林紀之 2008
\item 公開番号2008-097341 改善活動支援プログラム、方法及び装置 有馬淳,\underline{鵜飼孝典},青山浩二,小林紀之 2008
\item 公開番号2008-0027927  振り返りデータ処理方法及び装置 \underline{鵜飼孝典} 2008
\item 出願番号2006-204277  振り返りデータ処理方法及び装置 \underline{鵜飼孝典} 2006
\item 公開番号2004-133742 作業支援方法、作業支援プログラムおよび作業支援システム \underline{鵜飼孝典},片山佳則,古川淳子 2004
\item 登録番号7424483 配置情報推薦装置、その方法、プログラム \underline{鵜飼孝典},三末和男 2003
\item 登録番号4118580 配置情報推薦装置、その方法、プログラム \underline{鵜飼孝典},三末和男 2003
\item 登録番号4018919  ディレクトリ流通管理装置及び方法 \underline{鵜飼孝典},三末和男 2002
\item 登録番号6873983  キーワードの利用頻度予測装置 \underline{鵜飼孝典},三末和男 2002
\item 公開番号2002-351897  情報利用頻度予測プログラム、情報利用頻度予測装置および情報利用頻度予測方法 \underline{鵜飼孝典},三末和男 2002
\item 公開番号2001-117808 文書ディレクトリの分散管理システムおよびその取得方法並びにその方法をコンピュータに実行させるプログラムを記録したコンピュータ読み取り可能な記録媒体 \underline{鵜飼孝典} 2001
\item 公開番号H10-098452 通信文受信装置 蓬莱尚幸,\underline{鵜飼孝典} 1999
\item 公開番号H10-177342 文字置き換え暗号解析装置，その方法およびそのプログラム記録媒体 \underline{鵜飼孝典},蓬莱尚幸 1999
\item 出願番号H08-270461  文字置き換え暗号解析装置，その方法およびそのプログラム記録媒体 \underline{鵜飼孝典},蓬莱尚幸 1997
\item 公開番号H08-202726 分散データベースシステム及びそのデータ転送方法 \underline{鵜飼孝典} 1997
\item 公開番号H08-030558 計算機システムにおける負荷分散方法及びそれを利用した計算機システム \underline{鵜飼孝典},上田晴康 1996
\item 公開番号H07-334356  ソフトウェアの開発過程記述支援方法および装置  \underline{鵜飼孝典} 1995
\end{enumerate}
\item その他
\begin{enumerate}
\item プレスリリース
\begin{enumerate}
\item 富士通がヨーロッパで新しいAIおよびデータ分析のフレームワークSholarkを発表、販売開始 2018.11.06 \url{https://www.fujitsu.com/emeia/about/resources/news/press-releases/2018/emeai-20181106-new-fujitsu-framework-sholark-accelerates.html}
\end{enumerate}
\item 学会誌特集編集
\begin{enumerate}
\item 情報処理2013年10月号別刷「《特集》ユーザスタディのフロンティア」\\
編集者 \underline{鵜飼孝典}, 田村 大
\end{enumerate}
\end{enumerate}
\end{enumerate}

\newpage
\setcounter{section}{4}
\section{その他の活動}
\begin{enumerate}
\item 学会、社会、国際貢献
\begin{itemize}
\item 2024年4月～現在 人工知能学会セマンティックウェブとオントロジー研究会主査
\item 2024年 CIKM(International Conference on Information and Knowledge Management) 2024 Program Committee of Short Paper
\item 2022年11月～2023年12月 The 12th International Joint Conference on Knowledge Graphs (IJCKG 2023) Sponsorship Chair
\item 2023年 CIKM(International Conference on Information and Knowledge Management) 2023 Program Committee of Resource Track
\item 2023年 KGR4XAI (International Workshop on Knowledge Graph Reasoning for Explainable Artificial Intelligence) 2023 プログラム委員
\item 2021年1月～現在 ヒューマンインタフェース学会ユーザエクスペリエンス及びサービスデザイン専門研究委員会専門委員
\item ISWC(International Semantic Web Conference) 2019,2020,2021,2022,2023,2024,2025,2026 プログラム委員
\item KGSWC(Knowledge Graphs and Semantic Web Conference) 2019,2020,2021,2022,2023,2024,2025 プログラム委員
\item 2020年4月～2024年3月 人工知能学会セマンティックウェブとオントロジー研究会幹事
\item 人工知能学会合同研究会2020実行委員
\item JIST(Joint International Semantic Technology Conference) 2018 組織委員
\item 2018年1月～2020年12月 ヒューマンインタフェース学会ユーザエクスペリエンス及びサービスデザイン専門研究委員会幹事
\item 2016年1月～2017年12月 ヒューマンインタフェース学会ユーザエクスペリエンス及びサービスデザイン専門研究委員会専門委員
\item EPIC(Ethnographic Praxis in Industry Conference) 2010 開催地委員
\item CollabTech(International Conference on Collaboration Technologies) 2009,2012 プログラム委員
\item KICSS(International Conference on Knowledge, Information and Creativity Support Systems) 2008 プログラム委員
\item CollabTech(International Conference on Collaboration Technologies) 2006,2008 実行委員
\item 2005年4月～2008年3月 情報処理学会グループウェアとネットワークサービス研究会幹事
\item ICDCS(International Conference on Distributed Computing Systems) 2004 実行委員
\item 2003年人工知能学会全国大会プログラム委員
\item 2001年4月～2005年3月 情報処理学会グループウェアとネットワークサービス研究会運営委員
\item 1995年4月～1996年7月 情報処理学会 情報規格調査会 FDT-SWG 委員
\end{itemize}
\item 受賞
\begin{itemize}
\item 2025年11月 Seventh International Knowledge Graph and Semantic Web Conference Best Research Paper (Comparison of Metadata Representation Models for Knowledge Graph Embeddings) 受賞
\item 2023年6月 人工知能学会 2022年度人工知能学会研究会優秀賞 (イベント中心ナレッジグラフ埋め込みにおけるメタデータ表現モデルの分析) 受賞
\item 2022年12月 Linked Open Data チャレンジ Japan 2022 データ作成部門優秀賞 (VirtualHome2KGデータセット ― 家庭内の日常生活行動のシミュレーション動画とナレッジグラフ ―) 受賞
\item 2022年11月 人工知能学会 全国大会優秀賞 (説明生成のための知識グラフ構築ガイドラインの考察 ― ナレッジグラフ推論チャレンジを例にして ‒) 受賞
\item 2020年 人工知能学会ナレッジグラフ推論チャレンジベストツール賞 受賞
\item 2018年 人工知能学会ナレッジグラフ推論チャレンジベストリソース賞 受賞
\item 2018年 人工知能学会ナレッジグラフ推論チャレンジ優秀賞 受賞
\item 2017年 内閣府主催RESASアプリコンテスト ソフトバンクテクノロジー賞 受賞
\item 2017年 LODChallenge2016 審査員特別賞 受賞
\item 2011年 情報処理学会グループウェアとネットワーク研究会優秀発表賞 受賞
\item 2009年 情報処理学会 平成21年度山下記念研究賞 受賞
\item 2008年 情報処理学会グループウェアとネットワーク研究会優秀発表賞 受賞
\item 2001年 人工知能学会 全国大会優秀賞 (特定ドメイン向けWebディレクトリのためのWebリンク解析) 受賞
\end{itemize}
\end{enumerate}
\section{教育活動及び研究活動に係る今後の計画・抱負等}

\begin{enumerate}
\item 【「教育（人材育成）の基本方針」，その実現に向かっての抱負，講義・訓練等（科目等），研
究指導，学修指導において取り組む内容，目標，工夫について】\\
\textbf{教育（人材育成）の基本方針}について、自分の持つ知識を広く社会と共有し、自分の価値を社会に問える人材を育てることと考えています。
この方針の下で研究成果は学会などで積極的に発表を行うことを推奨します。
北陸先端科学技術大学院大学の学生は原則大学の教育を修了してきていますので、できるだけ早い段階から積極的に研究会などで
発表し、多くの人に理解できるように説明し、建設的な議論を行い、自身の研究を深めていくことが重要だと考えています。
目標としては、研究室で担当する学生については、過程の途中で1回、最終成果について1回、最終成果については
できれば英語で海外での発表を行わせたいと考えています。
そのためには普段の講義の中で、他人が理解できるように説明し、建設的な議論を行う技術を学生に伝えることが必要だと考えます。
社内や学会での経験から特に建設的な議論を引き出すような質問力を向上させるような教育が必要だと考えます。
講義の中で相互に質問し、その質問を振り返ることを習慣にしたいと考えています。
この方法は、私が英国ケンブリッジに駐在していた時に親しくしていた教師が教室で行っていた方法で
教室を見学させてもらってたいへん効果的だと感じたものです。

また、制作したソフトウェアは、OSSとして公開して、多くの人の協力を得て、より高品質で、
より多くの人に使われるものにしていくことを推奨したいと考えています。
あるいは、オンラインでの販売を行わせることも推奨したいと考えています。
研究成果物は研究室や大学にとどめるのではなく、有償、無償での公開を原則にしたいと考えています。
成果物の作成については、すべてを自分で作成する必要はなく、むしろ既存のソフトウェアに新しい機能を追加する、
あるいは機能を変更するほうが望ましいと考えます。すでに存在する機能は利用し、自分の独自の機能に注力するのが
望ましいと考えます。これは既存研究の上に自分のアイディアを積み上げる科学、技術の研究と同じ考え方であると思います。
既存のソフトウェアに新しい機能を追加する場合、最終的には元のソフトウェアにその機能が取り込まれるように
指導していきたいと思います。自分が考えた機能が取り込まれるためには、その機能の有効性を使用例などを持って
説明し、多くの開発者に共感を持ってもらうことが必要です。これは論文が採録される過程と同じであり、論文執筆と
同じような価値があると考えます。この点において、私がいくつかのOSSで行ってきた経験をもとにした指導ができると
考えています。
また、その制作物を持って、学会などで開催されているチャレンジ（コンペティション）への応募、データソン、ハッカソンへの
参加も推奨したいと思います。これも自分の制作物の価値を社会に問うという点でとても有効な手段であると考えます。
有償、無償いずれの場合でも無責任に公開するだけではなく、公開後のサポートの方針を明示することが必須だと考えています。
自分ではサポートできないため、無保証、サポートなしの場合もあると想定されますが、それを明らかにして
ユーザに伝えることが責任ある公開の仕方であると考え、そのように指導していきたいと考えます。

自分の論文や制作物の公開について、「良い評価を得られないから公開したくない」という反応を得たことがあります。
私は、多くの人からフィードバックをもらい、自分の作ったものをより良いものにしていくために公開すると考えています。
そのように学生に伝え、指導していきたいと考えています。これを実現するために講義においても積極的に学外の方との
コミュニケーションを行うことができるようにしたいと考えています。例えば、提出させるレポートを発表形式にし、
知り合いお研究者、技術者に協力をお願いし、ビデオ配信を行い、直接フィードバックを得るようにしたいと考えています。

\item 採用後１０年間の研究計画（国際連携・産学連携活動を含む），抱負，達成目標，及びそれ
らの見通しについて\\
2019年から人工知能学会のセマンティックウェブとオントロジー研究会のメンバーとともにナレッジグラフ推論チャレンジを主催し、
納得できる説明を行うAIの研究、開発の推進を行っています。これまで3回国内でイベントを行い、2021年から国際会議のイベントとして
実施すべく準備を進めています。このイベントでは、シャーロックホームズの小説を用いて、シャーロックホームズと同じような
事件の解決を行うことを目標としています。推理小説を解く課題設定でのイベントは、国際イベント3回くらいでめどをつける計画で、
次の段階で、扱う課題をより実践的な社会課題に向けるべく課題の検討を開始しています。
説明性は社会的に現在AIに求められている重要な課題であるため、この活動は引き続き行っていきたいと
考えています。

また、今年度から産業総合研究所人工知能研究センターと共同で5年計画で背景知識を用いて現場に起こる課題を予測するAIの研究を始めました。
このプロジェクトはNEDOの委託事業を受託して行っています。
具体的には介護の場面を適用例として用いて、そこで起こるヒヤリハットを予測する技術の研究を行っています。
医薬、介護をはじめとして様々な分野でオントロジーとして多くの知識が体系化され一般的に利用できるようになってきています。
このプロジェクトでは機械学習による予測の精度を上げるためにそれらの背景知識を有効に利用する技術の研究開発が目標になっています。
このプロジェクトについても引き続き参画したいと考えています。

さらに現在の勤務先では、知識発見に改めて取り組むことになりました。過去に行ってきたデータマイニングに基づいてパターンを
抽出する知識発見の研究、近年の説明可能AIの研究から、実務的なデータに技術を適用すると「専門家なら当たり前に知っている
知識しか示すことができない」あるいは「実務の役に立たない些細な知識しか示すことができない」という反省が得られました。
そのため、あらためて人間はどのように知識の発見をしてきたかという問いに立ち戻り、新しい知識の発見を行う技術の研究に
着手しています。観察事象の抽象化による新たな仮説の生成、背景知識と事象の矛盾の発見による新たな仮説の生成などに
取り組み始めています。この研究については、医療、マテリアル、マーケティングなどの応用領域や
説明可能AI、知識の検証の技術に関して国内、海外の複数の研究機関、大学と組んで行っています。
この分野に関しては、採用後も共同研究として引き続き取り組んでいきたいと考えています。
企業で研究を行っているときは、医療、マテリアル、マーケティングなど実利的な応用領域に取り組んでいましたが
大学で採用された場合には、文学、考古学、音楽なども応用領域として加えたいと考えています。
近年博物館、図書館などで芸術などのデータの公開が進んでいますが、必ずしも十分に活用できてはいないと考えます。
知識発見の技術を応用することで、これらの分野における研究を支援することもあるいは、新しい芸術作品を
作り出すことも可能になると考えています。

今回採用された場合、新しく次の研究を立ち上げて、実現することを10年後の目標にしたいと考えています。
産業総合研究所との共同プロジェクトを発展させて、組織経営、マーケティングに関する背景知識から
組織の知識経営を支援する技術の研究です。これには背景知識としてのオントロジーの利用技術に関して様々な論文、著作が
あるスペインUniversidad Politecnica de MadridのOscar Corcho教授の協力を仰ぎたいと考えています。
Corcho教授には、スペインでの精神医療のAIを研究している際に大変有効な助言をいただきました。
近年デジタルトランスフォーメーションという言葉が生まれ、データやAIを利用してより良い経営を
行うことが進められています。しかしながら直接的に経営方針を示唆する有効な技術はいまだ実現できていないと
考えます。組織の運営方針、施策を背景知識や経営データから仮説として生成し、それを実施した場合の数年後の
組織の状態の予測を行い、その予測の根拠を示すことで経営者がより適切な判断ができるようになると考えています。
これは、先に述べた3つの研究が順調に成果を出せれば、その組み合わせにより実現できると考えています。


\item 学生獲得，就職支援活動など大学の業務及び管理運営に関する抱負についての具体的見通しについて\\
現在の勤務先では、2010年からリクルータとして活動し、毎年数名との面接をおこなってきました。また今回の応募について
勤務先の人事部本部長から推薦をもらいました。コロナ感染症の広がりとタイミングを合わせて
テレワークが広がり、働き方改革が進んできており、企業が求める人物像も急速に変化してきております。
入社し年月が経つと上司の推薦により管理職になるということはなくなりつつあり、キャリアパスを自分で設計し、
必要に応じて自ら勉強したり、資格を取ったり、転職という形ばかりでなく社内においても自らポジションを求めることが求められる
ようになってきております。キャリア形成のために副業や兼業などの選択肢も考えられるようになってきております。
この変化に対応するため自分の経験を交えて学生に指導していきたいと思います。

また、英国駐在時は、15の企業で構成されるコンソーシアムで研究をするとともに会社の代表として意見を伝えることも
おこなってきました。その他にも海外企業との共同研究のプロジェクトに複数参加してきました。
特に米国PARC社との共同プロジェクトでは、PARC社の文化人類学のチームとの協業を経験し、
社内にフィールドイノベーターという新しい職種を作り出し、その教育プログラム作成にも貢献しました。
これらの経験から立場、考え方の異なる先生方を理解しつつ、お互いに協力して
大学としての強さを社会に出していけると考えています。
\end{enumerate}
\newpage
\section{応募者本人の連絡先}
\begin{description}[style=sameline]
\item [住所]　　〒182-0021　東京都調布市調布ヶ丘3-22-5
\item [Tel]　 　(042)485-3396
\item [E-Mail]　\href{mailto:takanori.ugai@gmail.com}{takanori.ugai@gmail.com}
\end{description}
\newpage
\section{意見をうかがえる3人の氏名・所属と連絡先(Eメールアドレス)}
\begin{itemize}
\item (富士通株式会社 富士通研究所 人工知能研究所 数理発見プロジェクト プロジェクトディレクタ) (\href{mailto:ohori.kotaro@fujitsu.com}{ohori.kotaro@fujitsu.com})：職場上司
\item 三末 和男(筑波大学システム情報系教授) (\href{mailto:misue@cs.tsukuba.ac.jp}{misue@cs.tsukuba.ac.jp})：元職場の上司であり、筑波大学における非常勤講師のホスト
\item 古崎 晃司(大阪電気通信大学 情報通信工学部 情報工学科教授) (\href{mailto:kozaki@osakac.ac.jp}{kozaki@osakac.ac.jp})：共同で取り組んでいるナレッジグラフ推論チャレンジ企画運営のリーダー
\end{itemize}
\end{document}
